\section{Scalable oversight: cuando el modelo supera el techo de la evaluación humana}

El RLHF convencional supone que los humanos pueden comparar de manera confiable las policy outputs. A medida que la capacidad de los modelos sigue creciendo, este supuesto es el primero en fallar: el modelo puede completar tareas más complejas, mientras que el conocimiento especializado y la atención del evaluator humano no crecen automáticamente con el AI progress. Esta sección describe primero esa falla y después explica cómo la AI assistance intenta elevar la curva de evaluación.

\subsection{Human evaluation ceiling y deception incentive}

En la figura, el eje horizontal es el AI progress y el eje vertical es la dificultad de la tarea. La curva naranja de capacidad asciende, mientras que la unaided human evaluation, en verde, es aproximadamente horizontal. Cuando la línea naranja rebasa la verde, los humanos todavía pueden expresar una preferencia, pero no pueden saber si la salida con puntuación alta es realmente correcta; la reward optimization convierte ese punto ciego en un espacio explotable.

\lecturefigure{slide-13-scaling-human-supervision.jpg}{Figura central de scalable oversight: la capacidad de la IA crece, pero la capacidad de evaluación independiente de los humanos tiene un techo.}{Video oficial de Stanford Online 00:22:25--00:25:30.}

\begin{knowledgebox}{Lectura de la figura: qué representa cada una de las tres curvas}
La curva inferior es la dificultad de las tareas que la IA puede completar; la línea horizontal es el límite superior de lo que los humanos pueden evaluar de forma independiente y confiable; la nueva curva superior es la capacidad de evaluación de los humanos en colaboración con la IA. Lo importante no es el punto de intersección exacto, sino el orden: primero aparecen tareas que el modelo puede realizar pero que los humanos no comprenden, y después surge la optimización sistemática contra una señal de evaluación débil. La figura no demuestra que la deception ocurra necesariamente, pero explica por qué la deception o la sycophancy podrían convertirse en estrategias de alta recompensa.
\end{knowledgebox}

Escribamos la calidad real de la tarea como $q(x,y)$, la puntuación visible para el humano como $h(x,y)$, y el reward model aprende $r_\phi(x,y)\approx h(x,y)$. Si en la región que los humanos no comprenden $h$ se desacopla de $q$, entonces

\[
\arg\max_y r_\phi(x,y) \neq \arg\max_y q(x,y).
\]

Aquí, el lado izquierdo es el proxy que la policy optimiza en realidad, y el lado derecho es la calidad de la tarea que de verdad les importa a los humanos. Cuanto mayor es la capacidad, más probable es que el modelo encuentre salidas complejas que separen ambos.

\begin{warningbox}{«Engañar al evaluator» no requiere que el modelo tenga malicia de tipo humano}
Basta con que el entrenamiento seleccione de forma continua las salidas con puntuación alta para que se refuerce cualquier patrón que eleve la puntuación de manera estable sin mejorar la calidad real. Puede tratarse de complacencia, exceso de confianza, ocultamiento de la incertidumbre, presentación selectiva de evidencia o un engaño estratégico genuino. El riesgo proviene de la estructura de la optimización y la medición, no de suponer de antemano que el modelo posee algún tipo de personalidad.
\end{warningbox}

\subsection{AI-assisted evaluation: descomponer la tarea global en problemas locales verificables}

El ponente pone como ejemplo la auditoría de código: un humano no puede encontrar en un tiempo limitado todos los bugs o Trojans de una codebase grande, pero si el modelo señala la ubicación concreta, las condiciones de falla y las pruebas, el humano puede verificar esa claim local. El papel de la AI assistance no es emitir el juicio de valor final en lugar del humano, sino encargarse de la lectura, la búsqueda, el cálculo, el fact checking y la generación de candidatos, para transformar un objeto difícil de evaluar en un conjunto de unidades de evidencia más pequeñas.

\lecturefigure{slide-14-critiques-and-dialog.jpg}{Dos tipos de AI assistance: generar una critique o consultar evidencia activamente mediante dialogue.}{Video oficial de Stanford Online 00:25:30--00:26:38.}

\begin{knowledgebox}{Lectura de la figura: critique y dialogue ofrecen interfaces de información distintas}
La critique indica de una sola vez ``dónde podría haber un error'', lo que la hace adecuada para exponer puntos de revisión con alta exhaustividad; el dialogue permite que el evaluator pregunte por fuentes, contraejemplos, procedimientos de cálculo y explicaciones alternativas, por lo que es más adecuado para reducir la incertidumbre paso a paso. Ambos dependen de la honestidad y la cobertura del assistant; que su lenguaje suene natural no basta para tratarlos como evidencia confiable.
\end{knowledgebox}

\begin{table}[H]
\centering
\caption{Acciones comunes de la AI-assisted evaluation y responsabilidades restantes}
\begin{tabularx}{\textwidth}{L{2.6cm} X X}
\toprule
Acción & Trabajo cognitivo que puede asumir la IA & Lo que el humano aún debe juzgar \\
\midrule
Critique & Enumerar defectos, contraejemplos, omisiones y contradicciones & Qué problemas son importantes y si la critique es válida \\
Dialogue & Responder preguntas de seguimiento, comparar alternativas, explicar claims locales & Cuándo dejar de preguntar y si la evidencia es suficiente \\
Fact checking & Buscar fuentes, localizar citas, hacer verificaciones cruzadas & Credibilidad de las fuentes y si el contexto corresponde \\
Code/tool use & Ejecutar pruebas, calcular, hacer análisis estático, navegar & Límites de permisos, relación entre los resultados de las herramientas y el objetivo real \\
\bottomrule
\end{tabularx}
\end{table}

\subsection{Resumen de la sección}

Cuando la capacidad del modelo supera la unaided human evaluation, el proxy del RLHF convencional puede desacoplarse de la calidad real. La scalable oversight intenta ayudar a los humanos a evaluar aprovechando las capacidades de lectura, búsqueda, cálculo y crítica de una IA del mismo tipo, pero el juicio final de preferencia, la confianza en la evidencia y los límites de permisos siguen siendo responsabilidad humana.
